\lecture{2.7}{Marginal and Conditional Causal Effects}
\label{lec:2-7}

\subsection{Marginal ATE}

The marginal average treatment effect summarizes the effect in the entire target population:
\[
\ATE=\E(\Yone-\Yzero).
\]
It averages over all sources of variation in the population, including baseline covariates. In the early imaging example, the marginal ATE answers:
\begin{quote}
What is the average difference in disability if everyone in the target population received early imaging versus if no one received early imaging?
\end{quote}

The marginal ATE is often the primary public health or policy estimand because it describes the average impact of a population-level intervention.

\subsection{Conditional ATE}

A conditional average treatment effect restricts the comparison to a subgroup defined by covariates.

\begin{definition}[Conditional average treatment effect]
For covariates $W$, the \emph{conditional average treatment effect} at covariate value $w$ is
\[
\CATE(w)=\E(\Yone-\Yzero\given W=w).
\]
\end{definition}

If $W$ is baseline pain severity, then $\CATE(w)$ is the average effect of early imaging among patients with baseline pain severity $w$.

\subsection{Subgroup-Specific Causal Effects}

Subgroup-specific causal effects are conditional effects for scientifically meaningful groups. If $W$ is categorical, such as low, moderate, and severe baseline pain, then
\[
\E(\Yone-\Yzero\given W=\text{severe})
\]
is the ATE among patients with severe baseline pain.

The marginal ATE is a weighted average of conditional ATEs:
\[
\ATE
=\E\{\E(\Yone-\Yzero\given W)\}.
\]
For discrete $W$,
\[
\ATE
=
\sum_w \E(\Yone-\Yzero\given W=w)\Prob(W=w).
\]
This is the tower rule applied to the individual-level contrast $\Yone-\Yzero$.

\subsection{Causal Effect Modification}

\begin{definition}[Causal effect modification]
There is \emph{causal effect modification} by $W$ on the additive scale if
\[
\E(\Yone-\Yzero\given W=w)
\]
varies with $w$.
\end{definition}

For example, early imaging might have little average effect among patients with mild baseline pain but a different effect among patients with severe neurological symptoms. If the conditional causal effect differs across baseline severity levels, severity modifies the causal effect.

\begin{example}[A finite population]
Suppose a population of ten individuals has potential outcomes shown below.
\[
\begin{array}{c|ccc|c}
\text{ID} & W & \Yone & \Yzero & \Yone-\Yzero\\
\hline
1 & 2 & 2 & 3 & -1\\
2 & 1 & 10 & 5 & 5\\
3 & 1 & 8 & 7 & 1\\
4 & 1 & 6 & 3 & 3\\
5 & 0 & 2 & 0 & 2\\
6 & 2 & 3 & 6 & -3\\
7 & 0 & 4 & 0 & 4\\
8 & 1 & 0 & 5 & -5\\
9 & 2 & 7 & 0 & 7\\
10 & 0 & 6 & 9 & -3
\end{array}
\]
The marginal ATE is the average of the ten individual contrasts:
\[
\ATE=\frac{-1+5+1+3+2-3+4-5+7-3}{10}=1.
\]
The conditional effects are
\[
\CATE(0)=\frac{2+4-3}{3}=1,
\]
\[
\CATE(1)=\frac{5+1+3-5}{4}=1,
\]
\[
\CATE(2)=\frac{-1-3+7}{3}=1.
\]
In this example the marginal and subgroup effects are all equal on the additive scale, so there is no additive effect modification by $W$.
\end{example}

\subsection{Difference between Statistical Interaction and Causal Effect Heterogeneity}

Statistical interaction refers to variation in an observed-data association across levels of a covariate. For example, a regression model might show that
\[
\E(Y\given A=1,W=w)-\E(Y\given A=0,W=w)
\]
varies with $w$.

Causal effect heterogeneity refers to variation in a counterfactual contrast:
\[
\E(\Yone-\Yzero\given W=w).
\]
These are not the same object. Observed statistical interactions can arise from confounding, selection, measurement, or model scale. Conversely, causal effect heterogeneity can exist even when a particular regression model fails to show an interaction.

\begin{keyidea}
Effect modification is scale-specific and causal. It concerns whether a causal contrast, such as a risk difference or mean difference, varies across subgroups.
\end{keyidea}
